\section{Experimental Setup}
\label{sec:experimental}

\subsection{Dataset}
\label{sec:experimental:dataset}

We use the full official release of the Explainable Detection of Online
Sexism (EDOS) dataset from SemEval-2023 Task 10 \cite{kirk2023semeval}:
20,000 English-language social media posts labeled sexist / not sexist,
with every sexist post further annotated into one of four categories
(threats/plans to harm/incitement, derogation, animosity, and prejudiced
discussions). We use the dataset's own official train/dev/test partition
(14,000/2,000/4,000), fetched from a pinned upstream commit and verified
by SHA-256 checksum against the exact release used throughout this study,
so an upstream force-push cannot silently change what was evaluated.
Class balance is stable at roughly 24\% sexist / 76\% not sexist across
all three splits (train: 3,398 sexist / 10,602 not sexist; dev: 486 / 1,514;
test: 970 / 3,030, the exact figures underlying
Sec.~\ref{sec:results:overview}'s base rate).

The training split supplies both the few-shot demonstration pool
(Sec.~\ref{sec:methodology:prompts}) and the fine-tuning data for all
three baselines (Sec.~\ref{sec:methodology:baselines}); the dev split is
used only for the fine-tuned baselines' epoch selection
(Sec.~\ref{sec:methodology:baselines}). All reported metrics (prompting
grid, baselines, and sensitivity sweeps alike) are computed on the same
held-out test split. For the subtype-level error analysis
(Sec.~\ref{sec:results:subtype}), the test split's 970 sexist examples
break down as: derogation (454), animosity (333), prejudiced discussions
(94), and threats/plans to harm/incitement (89).

\subsection{Inference Precision}
\label{sec:experimental:precision}

Generation uses bf16 throughout, with no additional post-training
quantization: calibration showed 4-bit NF4 quantization and bf16 cost
effectively the same wall-clock time per example (30.28s vs.\ 30.47s), so
bf16 was kept as the default given ample VRAM headroom across the
roster, rather than paying the complexity of a quantized inference path
for no measured throughput benefit. This is separate from the QLoRA
baseline (Sec.~\ref{sec:methodology:baselines}), which always quantizes
its frozen base weights to 4-bit NF4 as standard QLoRA practice; there,
quantization is what makes fine-tuning a full-size 7--9B model on a
single GPU feasible at all, an entirely different concern from
generation-time speed.

\subsection{Generation and Decoding}
\label{sec:experimental:generation}

All generation is greedy (\texttt{do\_sample=False}, \texttt{top\_p=1.0},
no repetition penalty) under a fixed global seed (42) throughout.
Non-chain-of-thought variants use a 2-token generation budget, matching
their forced trailing \texttt{ANSWER:} cue
(Sec.~\ref{sec:methodology:prompts}); chain-of-thought variants use a
220-token budget, calibrated for the five non-reasoning models' brief
prose reasoning. qwen3 is the exception: its native thinking mode can
reason for far longer than 220 tokens before reaching a decision, so its
chain-of-thought units use a 1,024-token budget instead, directly
reflected in its substantially higher per-example latency
(Sec.~\ref{sec:results:cost}). Generation batch size is 32 for five of the
six models, calibrated to stay within available GPU memory (zero
out-of-memory failures observed at that size); qwen3 uses batch size
128, calibrated the same way for its longer generations
(fail\_ratio 0.008 at that batch size).

\subsection{Evaluation Metrics}
\label{sec:experimental:metrics}

We report accuracy and binary precision/recall/F1 with the sexist class as
positive, computed identically for every (model, variant) unit, baseline,
and sweep condition. A response is parsed by locating its \emph{last}
occurrence of an answer marker (\texttt{ANSWER:}, or a residual
chat-template artifact) and reading the YES/NO token that follows; this
single rule handles both the prompt-inserted marker of non-chain-of-thought
variants and the model-produced marker at the end of a chain-of-thought
response without special-casing either. A response that never reaches an
unambiguous YES/NO decision (including the deliberately unresolved case
of both words appearing with no leading marker) is a parse failure; the
fraction of failures is reported per unit as \texttt{fail\_ratio} and, in
the headline metrics used throughout this paper, failures default to a
``not sexist'' prediction. This default is a real cost only when
\texttt{fail\_ratio} is non-trivial, which is itself informative
(Sec.~\ref{sec:results:llama-refusal}).
